% Part:sets-functions-relations
% Chapter: size-of-sets
% Section: pairing

\documentclass[../../../include/open-logic-section]{subfiles}

\begin{document}

\olfileid{sfr}{siz}{pai}
\olsection{جوړوونکې تابعې او کوډونه}

\begin{explain}
د کانتور د زيګزاګ طريقه د $\Nat^n$ د شمېر وړ والے په انځوريز
ډول څرګندوي۔ خو راځئ پر هغه جدول تمرکز وکړو چې $\Nat^2$ ښيي۔
په جدول کښې د زيګزاګ کرښه تعقيبول او مقامونه شمېرل راښيي چې
$\tuple{1,2}$ له عدد~$7$ سره تړل شوې ده۔ خو ښه به وي که دا
په لا نېغه توګه حسابولے شو۔ يعنې ښه به وي چې د زيګزاګ شمېرنې
\emph{معکوس} تابع $g\colon \Nat^2 \to \Nat$ په لاس کښې ولرو،
داسې چې:
\[
g(\tuple{0,0}) = 0, \;
g(\tuple{0,1}) = 1, \;
g(\tuple{1,0}) = 2, \; \dots,  
g(\tuple{1,2}) = 7, \; \dots
\]
په دې سره په دقيق ډول حسابولے شو چې $\tuple{n, m}$ به زموږ
په شمېرنه کښې چېرته راشي۔

په حقيقت کښې، په دوو مشاهدو سره $g$ په نېغه توګه تعريفولے شو۔
لومړے: که د $n$م قطار او $m$م ستون قيمت~$v$ وي، نو د
$(n+1)$م قطار او $(m-1)$م ستون قيمت $v + 1$ دے۔ دويم: زموږ
د شمېرنې لومړے قطار له مثلثي عددونو جوړ دے چې له $0$، $1$،
$3$، $6$ او داسې نورو پيلېږي۔ $k$م مثلثي عدد د هغو طبيعي
عددونو جمع ده چې $< k$ دي، او په $k(k+1)/2$ حسابېږي۔ د دې دوو
مشاهدو په يوځاے کولو لاندې تابع په پام کښې ونيسئ:
\[
  g(n,m) = \frac{(n+m+1)(n+m)}{2} + n
\]
ډېر وخت د $g(\tuple{n, m})$ پر ځاے يوازې $g(n, m)$ ليکو، ځکه
چې سترګو ته اسانه دے۔ دا لومړے د $(n+m)^\text{th}$ مثلثي عدد
ټاکل او بيا $n$ ورزياتول درته وايي۔ او دا جدول په هماغه طريقه
ډکوي چې غواړو ئې۔ نو په ځانګړي ډول، جوړه $\tuple{1, 2}$ له
$\frac{4 \times 3}{2} + 1 = 7$ سره تړل کېږي۔

دا تابع $g$ د جوړو د يوه سټ د شمېرنې \emph{معکوس} ده۔ داسې
تابعې \emph{جوړوونکې تابعې} بلل کېږي۔
\end{explain}

\begin{defn}[جوړوونکې تابع]
  تابع $f\colon A \times B \to \Nat$ يوه حسابي
  \emph{جوړوونکې تابع} ده که $f$ يو پر يو وي۔ دا هم وايو چې
  $f$ د $A \times B$ \emph{کوډول} کوي، او $f(x,y)$ د
  $\tuple{x,y}$ \emph{کوډ} دے۔
\end{defn}

\begin{explain}
جوړوونکې تابعې د بېلګې په توګه د طبيعي عددونو د جوړو د
کوډولو دپاره کارولے شو؛ يا په بله وينا، د غړو هره \emph{جوړه}
په يوه \emph{واحد} عدد ښودلے شو۔ د جوړوونکې تابعې په معکوس
عدد بېرته لوستلے شو؛ يعنې معلومولے شو چې کومه جوړه ښيي۔
\end{explain}

\begin{prob}
د ټولو غير منفي ناطقو عددونو د سټ يوه شمېرنه ورکړئ۔
\end{prob}

\begin{prob}
وښيئ چې $\Rat$ !!{enumerable} دے۔ په ياد ولرئ چې هر ناطق عدد
د کسر $z/m$ په توګه ليکل کېدے شي، چې $z \in \Int$ او
$m \in \Nat^+$ وي۔
\end{prob}

\begin{prob}
د $\Bin^*$ يوه شمېرنه تعريف کړئ۔
\end{prob}

\begin{prob}
د منطق له پېژندنيز کورسه په ياد ولرئ چې د رښتياوالي هر ممکن
جدول د رښتياوالي يوه تابع څرګندوي۔ په بله وينا، د کوم~$k$
دپاره د رښتياوالي تابعې ټولې تابعې له $\Bin^k \to \Bin$ څخه
دي۔ ثابته کړئ چې د رښتياوالي د ټولو تابعو سټ د شمېر وړ دے۔
\end{prob}

\begin{prob}
وښيئ چې د يوه خپل سري نامتناهي !!{enumerable} سټ د ټولو
متناهي فرعي سټونو سټ !!{enumerable} دے۔
\end{prob}

\begin{prob}
د $\Nat$ يوه فرعي سټ ته هله \emph{د متناهي متمم لرونکے} ويل
کېږي چې متمم ئې په $\Nat$ کښې متناهي سټ وي؛ يعنې
$A \subseteq \Nat$ هله د متناهي متمم لرونکے دے چې
$\Nat\setminus A$ متناهي وي۔ \textbf{د ژباړې د نسخې سمون
(OLSIZ-002):} د انګرېزي سرچينې په تعريف کښې د متناهي فرعي سټ او
طبيعي عددونو د سټ تر منځ اړيکه پاتې وه؛ ورپسې هم‌ارز فارمول ښيي
چې مطلب په طبيعي عددونو کښې متناهي متمم دے۔ فرض کړئ $I$ هغه سټ دے چې
!!{element}s ئې په دقيقه توګه د $\Nat$ متناهي او د متناهي
متمم لرونکي فرعي سټونه دي۔ وښيئ چې $I$ !!{enumerable} دے۔
\end{prob}

\begin{prob}
وښيئ چې د !!{enumerable} سټونو !!{enumerable} اتحاد
!!{enumerable} دے۔ يعنې، هرکله چې $A_1$، $A_2$، \dots{} سټونه
وي او هر $A_i$ !!{enumerable} وي، نو د هغو ټولو اتحاد
$\bigcup_{i=1}^\infty A_i$ هم !!{enumerable} دے۔ [يادونه:
دا ګران دے!]
\end{prob}

\begin{prob}
فرض کړئ $f \colon A \times B \to \Nat$ يوه خپله سري جوړوونکې
تابع ده او د جوړو سټ ناتش دے۔ وښيئ چې د~$f$ معکوس پر خپلې
د تعريف ساحې دوه اړخيزه يو پر يو تابع ده چې د قيمتونو سټ ئې
$A \times B$ دے، او د هغې له ساحې څخه بهر د يوې ثابتې جوړې
په ورکولو د طبيعي عددونو پر ټولو غړو غځېدے شي۔
\textbf{د ژباړې د نسخې سمون (PSSIZ-001):} انګرېزي تمرين د
معکوس تابع د تعريف ساحه د اصلي تابعې له قيمتونو سټ څخه ټولو
طبيعي عددونو ته نۀ وه غځولې، او د تش جوړو سټ حالت ئې هم نۀ و بېل کړے۔
\end{prob}

\begin{prob}
داسې تابع وټاکئ چې $\Nat^3$ کوډ کړي۔
\end{prob}

\end{document}
